\documentclass[11pt]{article}
\usepackage[a4paper,margin=27mm]{geometry}
\usepackage{amsmath,amssymb,amsthm}
\usepackage[hidelinks]{hyperref}
\newtheorem{theorem}{Theorem}
\newtheorem{lemma}[theorem]{Lemma}
\newtheorem{corollary}[theorem]{Corollary}
\newcommand{\Z}{\mathbb Z}
\newcommand{\sgn}{\operatorname{sgn}}
\title{Periodic two-cycles for every symmetric median mask\\
\large An elementary projection and an explicit period bound}
\author{MathIsEvenEasier\\\small Prepared with OpenAI Codex (GPT-6 Astra)}
\date{6 October 2026 --- research draft}
\begin{document}
\maketitle
\begin{abstract}
We revisit the August 2026 archived antipodal-cut construction for
nontrivial centered symmetric median masks on $\Z^d$. Its convex-geometric
separation step can be replaced by a balanced-base encoding of integer
coordinates. This gives a complete elementary proof and an explicit bound
on a spatial period. The existence statement, explicit spatial-period bound, and improving-flip
bound are formalized in Lean 4. Literature priority remains unconfirmed.
\end{abstract}

\section{Statement and scope}
Let $d\ge1$ and let
\[
 W=\{0\}\cup\{\pm a_1,\ldots,\pm a_r\}\subset\Z^d,
 \qquad r\ge1,
\]
where the $a_i$ are nonzero and distinct up to sign. For
$x:\Z^d\to\{-1,1\}$, the synchronous binary median is
\[
 F_W(x)(v)=\sgn\left(x(v)+\sum_{i=1}^r
                 (x(v+a_i)+x(v-a_i))\right).
\]
There are $2r+1$ samples, so the sum never vanishes.

\begin{theorem}\label{thm:main}
Every such mask admits a coordinate-periodic configuration $x$ with
$F_W(x)=-x$. If every coordinate of every $a_i$ has absolute value at most
$M\ge1$, then a common coordinate period $N$ can be chosen with
\[
             2\le N\le (2M+1)^d-1.
\]
Thus $x,-x$ form a genuine temporal two-cycle.
\end{theorem}

The period $N$ need not be minimal. The theorem concerns existence on
some finite torus, or equivalently a periodic image on the infinite lattice.
It makes no assertion for every prescribed image size, asymmetric masks,
weighted masks, or fixed-zero, reflected, or replicated boundary rules.
It addresses the centered-symmetric, periodic subclass of Klette's
Problem 3~\cite{klette}, rather than the entire characterization question.

\section{The elementary projection}
\begin{lemma}\label{lem:encode}
Set $B=2M+1$ and
$\ell(v)=\sum_{j=0}^{d-1} B^j v_j$.
Then $\ell$ is injective on the integer box $[-M,M]^d$.
Consequently the positive numbers $|\ell(a_i)|$ are pairwise distinct.
\end{lemma}
\begin{proof}
For distinct $u,v$ in the box, let $k$ be the largest index for which
$u_k\ne v_k$. The leading term has absolute value at least $B^k$, whereas
\[
 \left|\sum_{j<k}B^j(u_j-v_j)\right|
 \le 2M\sum_{j<k}B^j=B^k-1.
\]
Therefore $\ell(u)-\ell(v)\ne0$. Apply this to $0,a_i,-a_i$ and
to pairs $a_i,\pm a_j$ to obtain the last assertion.
\end{proof}

Relabel and orient the representatives so that
\[
 p=\ell(a_r)=\max_i|\ell(a_i)|>0,
 \qquad q_i=\ell(a_i),\quad 0<|q_i|<p\quad(i<r).
\]
Also $p\le M\sum_{j<d}B^j=(B^d-1)/2$. No exposing vertex or
rational-separation theorem is needed.

\section{Antipodal cut and majority}
Work in $C=\Z/(2p\Z)$. Choose a coloring $y:C\to\{-1,1\}$ satisfying
\[
                 y(t+p)=-y(t)
\]
which maximizes the integer potential
\[
 \Phi(y)=\sum_{i<r}\sum_{t\in C}
                    \mathbf 1_{y(t)\ne y(t+q_i)}.
\]
The choices are nonempty and finite: one sign is free at each of the $p$
antipodal pairs. Edge labels $(i,t)$ are retained throughout.
For each $t$, let
\[
 c(t)=\sum_{i<r}\left(
  \mathbf 1_{y(t)\ne y(t+q_i)}+
  \mathbf 1_{y(t)\ne y(t-q_i)}\right).
\]

\begin{lemma}\label{lem:flip}
At such a maximum, $c(t)\ge r-1$ for every $t\in C$.
\end{lemma}
\begin{proof}
Translate by $p$: both endpoint signs reverse, so $c(t+p)=c(t)$.
No remaining edge is a loop or joins an antipodal pair, since
$0<|q_i|<p$. For a fixed label $i$, the incoming and outgoing edge at a
vertex are distinct, since $2q_i\not\equiv0\pmod{2p}$.
Thus exactly $4(r-1)$ labeled edges meet the pair $\{t,t+p\}$, of which
$2c(t)$ are cut. Flipping both signs preserves antipodality and toggles
exactly those edges. The potential change is
\[
       \Delta\Phi=4(r-1)-4c(t).
\]
Maximality makes this nonpositive, proving the claim.
\end{proof}

\begin{proof}[Proof of Theorem~\ref{thm:main}]
Define $x(v)=y(\ell(v)\bmod 2p)$. This has common coordinate period
$N=2p$. The two samples at $v\pm a_r$ both project to the antipode of
$\ell(v)$, so both disagree with $x(v)$. Lemma~\ref{lem:flip} supplies at
least $r-1$ further disagreements. Hence at least $r+1$ of the $2r+1$
samples disagree with the center, proving $F_W(x)=-x$.
The operator commutes with sign reversal, whence $F_W(-x)=x$.
Finally $x\ne-x$ and $2p\le B^d-1$.
\end{proof}

\section{Constructivity and finite images}
The maximizing coloring need not be found by exhaustive search. Begin with
any antipodal coloring and flip a pair with $c(t)<r-1$. Each flip increases
$\Phi$ by at least four; since $0\le\Phi\le2p(r-1)$, there are at most
\[
                   \left\lfloor\frac{p(r-1)}2\right\rfloor
\]
improving flips. This gives a finite constructive procedure. It is not a
polynomial-time claim in the binary input length: $p$ can be large compared
with the bit length of the coordinates.

The witness is represented by $2p$ signs and the linear map $\ell$; it need
not be expanded into an $(2p)^d$ array. On a torus of side $2p$, different
mask offsets may coincide, and all samples must still be counted. If one
requires distinct finite-image neighbors, choose any side length $L$ that
is a multiple of $2p$ and satisfies $L>2M$. The same witness then descends
to that torus with all mask offsets distinct.

For the simple mask $\{0,\pm1,\pm2\}$, the repeated binary word $0011$
maps to $1100$ and back. The theorem gives such an oscillating test input
for every allowed mask. Consequently an iteration loop which stops only
when two consecutive images agree is not guaranteed to terminate under
circular boundary conditions. Comparing with the image two steps earlier
also detects the cycles, although it does not by itself select a preferred
denoised image.

\section{Known upper bound on temporal periods}
The classical symmetric-threshold period-two phenomenon is background,
not a new result of this note; see, for example, the discussion and
references in~\cite{goles}. For completeness, on any fixed finite torus
let $A$ be the symmetric sample-count matrix and put $x_{t+1}=\sgn(Ax_t)$.
The integer energy $E_t=x_{t+1}^{T}Ax_t$ satisfies
\[
 E_t-E_{t-1}
 =\sum_v (x_{t+1}(v)-x_{t-1}(v))(Ax_t)(v)\ge0.
\]
Each changed coordinate contributes $2|(Ax_t)(v)|\ge2$, since the row
sample count is odd. Energy is bounded, so eventually $x_{t+1}=x_{t-1}$.
Thus every finite-image trajectory eventually has temporal period one or
two. Theorem~\ref{thm:main} supplies existence of period two on a suitable
torus for every nontrivial mask; the mask $\{0\}$ is the identity.

\section{Research status}
The antipodal-cut existence argument comes from the supplied research
archive dated 21 August 2026. This revision replaces its convex-geometric
projection step by Lemma~\ref{lem:encode}, extracts the explicit spatial
period and flip-count bounds, and clarifies the finite-image conventions.
These are refinements of that argument, not a claim to have discovered
the archived existence theorem today.

The new Lean development proves the existence statement end to end,
starting with integer-lattice offsets. It checks balanced-base injectivity,
the exact antipodal energy change, a finite minimizing coloring, the cyclic
median update, and the periodic lattice pullback. The endpoint
\texttt{MedianCycles.periodic\_median\_two\_cycle\_bounded} asserts the
spatial-period bound as well as
$F_W(F_W(x))=x$, $F_W(x)\ne x$, and $F_W(x)=-x$.
It is checked with Lean 4.34.0 and mathlib commit
\texttt{5ed2965256430c3649e86755f9576b54eca72435}.
Its only axiom dependencies are the standard \texttt{propext},
\texttt{Classical.choice}, and \texttt{Quot.sound}.
The separate endpoint
\texttt{MedianCycles.cyclic\_improving\_flip\_bound} bounds the length of
any actual sequence of antipodal flips satisfying $c(t)<r-1$. Its proof
relates the disagreement count to integer energy, proves a drop of at
least eight per improving flip, and bounds the total available energy
range. Both quantitative bounds in this note are therefore checked in Lean.
This is kernel checking by Lean,
not an independent-kernel audit. Source and public build records are at
\url{https://github.com/MathIsEvenEasier/median-cycles}.

The original question~\cite{klette} was checked directly. Modern work on
total cycles~\cite{goles} uses the same $F(x)=-x$ condition but does not
supply the arbitrary-mask lattice construction in its inspected results.
A recent complexity preprint~\cite{jain} was screened at abstract level
only; its full text was not retrieved. A bounded search across related
terminology and indexed GitHub, X, and VibeMathed results did not locate
this exact universal theorem. This is not evidence sufficient to establish
priority, nor an independent expert review.

\begin{thebibliography}{9}
\bibitem{klette} R. Klette (compiler), \emph{Open Problems in (Digital)
Geometry}, Dagstuhl, March 2004, Problem 3, slide 4.
\url{https://tc18.org/openProblems/Dagstuhl_2004.pdf}.
\bibitem{goles} E. Goles, P. Montealegre, M. Rios-Wilson and S. Sene,
\emph{Dynamical stability of threshold networks over undirected signed
graphs}, arXiv:2309.01854, version 3, especially Section 3 and Lemma 2.
\url{https://arxiv.org/html/2309.01854v3}.
\bibitem{jain} S. Jain, F. Stephan and H. Tang,
\emph{Strictly Unfriendly $k$-Partitions: Sharp Degree Thresholds and
ETH-Based Lower Bounds}, arXiv:2610.04440, 3 October 2026.
\url{https://arxiv.org/abs/2610.04440}.
\end{thebibliography}
\end{document}
